% Generated draft: requires mathematical and source review.
\documentclass[11pt]{article}
\usepackage[margin=25mm]{geometry}
\usepackage{fontspec}
\setmainfont{Latin Modern Roman}
\usepackage{amsmath,amssymb,mathtools,unicode-math}
\setmathfont{Latin Modern Math}
\usepackage{enumitem,needspace,xurl,hyperref}
\setlength{\parindent}{0pt}
\setlength{\parskip}{5pt}
\setlength{\emergencystretch}{3em}
\newcommand{\N}{\mathbb{N}}
\newcommand{\Z}{\mathbb{Z}}
\newcommand{\Q}{\mathbb{Q}}
\newcommand{\R}{\mathbb{R}}
\newcommand{\C}{\mathbb{C}}
\newcommand{\F}{\mathbb{F}}
\newcommand{\A}{\mathbb{A}}
\newcommand{\PP}{\mathbb{P}}
\newcommand{\E}{\mathbb{E}}
\newcommand{\eps}{\varepsilon}
\newcommand{\dd}{\,\mathrm{d}}
\begin{document}
\section*{UnsolvedMath ID: 30000223}
\textbf{Source problem code:} catalogue entry
\subsection{Definitions and mathematical statement}
Let $R=k[x_1,\ldots,x_n]$ be a polynomial ring over a field, and let $\mathfrak a\subset R$ be an ideal. The question is to characterize those $\mathfrak a$ for which there is one descending chain
\[
 \mathfrak a_1\supseteq\mathfrak a_2\supseteq\cdots
\]
cofinal with the powers of $\mathfrak a$, meaning
\[
 \forall s\ge1\ \exists t\ge1:\mathfrak a_t\subseteq\mathfrak a^s,
 \qquad
 \forall t\ge1\ \exists s\ge1:\mathfrak a^s\subseteq\mathfrak a_t,
\]
such that
\[
 \operatorname{Ext}^i_R(R/\mathfrak a_t,R)
 \longrightarrow\operatorname{Ext}^i_R(R/\mathfrak a_{t+1},R)
 \quad\text{is injective for every }i\ge0\text{ and }t\ge1.
\]
These natural maps arise from the quotient surjection $R/\mathfrak a_{t+1}\to R/\mathfrak a_t$ by contravariance of $\operatorname{Ext}$ in its first argument. The same chain must work simultaneously in all degrees; choosing a separate chain for each $i$ would be a different assertion. The source imposes neither homogeneity nor a characteristic restriction on the polynomial ring.

By cofinality the direct limit of this Ext system is $H^i_{\mathfrak a}(R)$. Thus the requirement is equivalently that every stage map into this local-cohomology direct limit is injective, for the chosen common chain and every degree. The question asks which ideals admit such a chain, not whether the ordinary powers themselves always have this property.

The literature assessment checked on 21 August 2026 reports positive cases for monomial ideals, Frobenius-power systems associated with $F$-pure quotients, and several characteristic-zero classes. It says a classification of all ideals remains unresolved. Exact hypotheses of the special-case theorems must be read in the supplied references; they are not promoted to an answer for every ideal. The partially-solved label is dated triage requiring independent verification.
\subsection{Short English statement}
Which polynomial-ring ideals admit a cofinal descending chain whose natural Ext transition maps are injective in every degree? One common chain must work for all degrees, and it may differ from the ordinary power chain.
\subsection{Sources}
\begin{itemize}
\item \emph{Kommutative Algebra} (2006), Oberwolfach Reports, DOI 10.4171/owr/2005/19, pp. 1073--1126, \url{https://doi.org/10.4171/owr/2005/19}; proposed year supplied as 2005.
\item Literature assessment checked 21 August 2026: original report, \url{https://doi.org/10.4171/owr/2005/19}; Singh--Walther, ``local cohomology and pure morphisms,'' \url{https://arxiv.org/abs/math/0701524}; ``local cohomology of Du Bois singularities,'' \url{https://arxiv.org/abs/1605.02755}.
\item Local UnsolvedMath registry, permanent ID $30000223$, exact source code \texttt{catalogue entry}, created 25 August 2026 and updated 16 September 2026. Supplied verification explicitly confirms the polynomial-ring hypothesis, subscripted ideal chain, power cofinality and direction of the Ext maps.
\end{itemize}
\end{document}
